% =====================================================================
%  Secao 2: Por que o Backtracking Existe?
% =====================================================================
\section{Por que o Backtracking Existe?}

Antes de estudar como o algoritmo funciona, é importante entender por que ele foi criado.

Considere o seguinte problema:

Uma senha possui quatro posições e cada posição pode conter um dos números 1, 2 ou 3. Quantas combinações existem?
\[
  3 \times 3 \times 3 \times 3 = 81 \text{ combinações.}
\]

Agora imagine uma senha de oito posições contendo dez possibilidades para cada posição.
\[
  10^{8} = 100{.}000{.}000 \text{ combinações.}
\]

Percebe-se rapidamente que o número de possibilidades cresce de forma extremamente acelerada.

Muitos problemas reais apresentam exatamente esse comportamento. Nesses casos, testar todas as combinações possíveis pode se tornar muito caro computacionalmente.

O Backtracking surge como uma forma organizada de explorar alternativas, descartando rapidamente caminhos que não possuem chance de gerar uma solução válida.

\figura{fig:explosao}%
  {A quantidade de combinações cresce de forma acelerada à medida que aumentam as posições e as opções por posição.}%
  {figuras/tikz/explosao}
